
\documentclass{article}
\usepackage{colortbl}
\usepackage{makecell}
\usepackage{multirow}
\usepackage{supertabular}

\begin{document}

\newcounter{utterance}

\centering \large Interaction Transcript for game `cladder', experiment `full\_v1.5\_default', episode 2927 with reflexion{-}llama{-}3.1{-}8b{-}instruct.
\vspace{24pt}

{ \footnotesize  \setcounter{utterance}{1}
\setlength{\tabcolsep}{0pt}
\begin{supertabular}{c@{$\;$}|p{.15\linewidth}@{}p{.15\linewidth}p{.15\linewidth}p{.15\linewidth}p{.15\linewidth}p{.15\linewidth}}
    \# & \multicolumn{2}{c}{Player} && \multicolumn{2}{c}{Game Master} \\
    \hline

    \theutterance \stepcounter{utterance}  
    & & & \multicolumn{4}{p{0.6\linewidth}}{
        \cellcolor[rgb]{0.9,0.9,0.9}{
            \makecell[{{p{\linewidth}}}]{
                \texttt{\tiny{[P1$\langle$GM]}}
                \texttt{You are an expert in causal inference. The following question is not a typical commonsense query, but rather a meticulously designed question created by a professor specializing in causal inference, intended to assess the students' mastery of the course content.} \\
\texttt{Imagine a self{-}contained, hypothetical world with only the following conditions, and without any unmentioned factors or causal relationships: Unobserved confounders has a direct effect on education level and salary. Proximity to a college has a direct effect on education level. Education level has a direct effect on salary. Unobserved confounders is unobserved. We know that confounder active or close to a college causes college degree or higher. confounder active and college degree or higher causes high salary. We observed the person lives close to a college and confounder active. Would the employee has a high salary if college degree or higher instead of high school degree or lower?} \\
\texttt{Start your answer with 'Yes' or 'No', followed by additional reasoning or evidence to support your explanation.} \\
            }
        }
    }
     \\ \\

    \theutterance \stepcounter{utterance}  
    & \multicolumn{4}{p{0.6\linewidth}}{
        \cellcolor[rgb]{0.561,0.737,0.561}{
            \makecell[{{p{\linewidth}}}]{
                \texttt{\tiny{[P1$\rangle$GM]}}
                \texttt{No} \\
\\ 
\texttt{In this hypothetical world, we have a complex causal relationship between the variables. To answer the question, we need to apply the principles of causal inference, specifically the concept of counterfactuals.} \\
\\ 
\texttt{Given the information:} \\
\\ 
\texttt{1. Unobserved confounders (U) has a direct effect on education level (E) and salary (S).} \\
\texttt{2. Proximity to a college (P) has a direct effect on education level (E).} \\
\texttt{3. Education level (E) has a direct effect on salary (S).} \\
\texttt{4. Unobserved confounders (U) is unobserved.} \\
\texttt{5. Confounder active (U=1) or close to a college (P=1) causes college degree or higher (E=1).} \\
\texttt{6. Confounder active (U=1) and college degree or higher (E=1) causes high salary (S=1).} \\
\texttt{7. We observed the person lives close to a college (P=1) and confounder active (U=1).} \\
\\ 
\texttt{We want to know if the employee would have a high salary (S=1) if they had a college degree or higher (E=1) instead of a high school degree or lower (E=0).} \\
\\ 
\texttt{To answer this question, we need to consider the counterfactual scenario where the person has a college degree or higher (E=1) but still lives close to a college (P=1) and has the confounder active (U=1).} \\
\\ 
\texttt{Since the confounder (U) is unobserved, we cannot directly estimate its effect on salary (S). However, we can use the information that confounder active (U=1) and college degree or higher (E=1) causes high salary (S=1) to inform our counterfactual analysis.} \\
\\ 
\texttt{In this scenario, the person already has a college degree or higher (E=1) and lives close to a college (P=1), so the effect of proximity to a college (P) on education level (E) is already accounted for. The only difference between the observed and counterfactual scenarios is the education level (E).} \\
\\ 
\texttt{To determine if the employee would have a high salary (S=1) if they had a college degree or higher (E=1) instead of a high school degree or lower (E=0), we need to consider the effect of education level (E) on salary (S) while holding the confounder (U) and proximity to a college (P) constant.} \\
\\ 
\texttt{Since the confounder (U) is unobserved, we cannot directly estimate its effect on salary (S). However, we can use the information that confounder active (U=1) and college degree or higher (E=1) causes high salary (S=1) to inform our counterfactual analysis.} \\
\\ 
\texttt{In this scenario, the person already has the confounder active (U=1) and lives close to a college (P} \\
            }
        }
    }
    & & \\ \\

    \theutterance \stepcounter{utterance}  
    & & & \multicolumn{2}{p{0.3\linewidth}}{
        \cellcolor[rgb]{0.9,0.9,0.9}{
            \makecell[{{p{\linewidth}}}]{
                \texttt{\tiny{[GM$|$GM]}}
                \texttt{No} \\
            }
        }
    }
    & & \\ \\

    \theutterance \stepcounter{utterance}  
    & & & \multicolumn{2}{p{0.3\linewidth}}{
        \cellcolor[rgb]{0.9,0.9,0.9}{
            \makecell[{{p{\linewidth}}}]{
                \texttt{\tiny{[GM$|$GM]}}
                \texttt{yes} \\
            }
        }
    }
    & & \\ \\

    \theutterance \stepcounter{utterance}  
    & & & \multicolumn{2}{p{0.3\linewidth}}{
        \cellcolor[rgb]{0.9,0.9,0.9}{
            \makecell[{{p{\linewidth}}}]{
                \texttt{\tiny{[GM$|$GM]}}
                \texttt{game\_result = LOSE} \\
            }
        }
    }
    & & \\ \\

    \theutterance \stepcounter{utterance}  
    & & \multicolumn{4}{p{0.6\linewidth}}{
        \cellcolor[rgb]{0.95,0.95,0.95}{
            \makecell[{{p{\linewidth}}}]{
                \texttt{\tiny{[P1$\leftrightarrow$P]}}
                \texttt{Post{-}episode Reflexion (not sent to GM):} \\
\texttt{New plan: } \\
\\ 
\texttt{In my previous attempt, I failed to provide a clear and concise answer to the question. I got bogged down in explaining the causal relationships and the concept of counterfactuals, but I didn't directly address the question. } \\
\\ 
\texttt{To improve my strategy, I will focus on clearly articulating the counterfactual scenario and the implications of the given information. I will also make sure to directly answer the question with a clear "Yes" or "No" and provide supporting evidence.} \\
\\ 
\texttt{In this case, I will re{-}evaluate the given information and consider the following:} \\
\\ 
\texttt{* The person lives close to a college (P=1) and has the confounder active (U=1), which causes a high salary (S=1) when combined with a college degree or higher (E=1).} \\
\texttt{* The effect of proximity to a college (P) on education level (E) is already accounted for, as the person already has a college degree or higher (E=1).} \\
\texttt{* The only difference between the observed and counterfactual scenarios is the education level (E).} \\
\\ 
\texttt{Given this, I will re{-}evaluate the causal relationships and consider the following:} \\
\\ 
\texttt{* If the person had a college degree or higher (E=1) instead of a high school degree or lower (E=0), would they still have the confounder active (U=1) and live close to a college (P=1)?} \\
\texttt{* If the answer is yes, then the effect of the confounder (U) and proximity to a college (P) on salary (S) would still be present, and the person would likely have a high salary (S=1).} \\
\\ 
\texttt{Therefore, my revised answer is:} \\
\\ 
\texttt{Yes, the employee would have a high salary (S=1) if they had a college degree or higher (E=1) instead of a high school degree or lower (E=0).} \\
            }
        }
    }
    & \\ \\

\end{supertabular}
}

\end{document}
